\section*{第 \ref{ch:g10:the-mole} 章——物质的量与摩尔质量}

\begin{solution}{exo:g10:the-mole:1}
$M(\ce{O2}) = 2 \times 16.0 = \qty{32.0}{g/mol}$;
$M(\ce{N2}) = 2 \times 14.0 = \qty{28.0}{g/mol}$;
$M(\ce{CH4}) = 12.0 + 4 \times 1.0 = \qty{16.0}{g/mol}$;
$M(\ce{NH3}) = 14.0 + 3 \times 1.0 = \qty{17.0}{g/mol}$;
$M(\ce{CO2}) = 12.0 + 2 \times 16.0 = \qty{44.0}{g/mol}$.
\end{solution}

\begin{solution}{exo:g10:the-mole:2}
水：$n = 36.0 / 18.0 = \qty{2.00}{mol}$。甲烷：
$n = 4.0 / 16.0 = \qty{0.25}{mol}$.
\end{solution}

\begin{solution}{exo:g10:the-mole:3}
$m(\ce{NaCl}) = 0.25 \times 58.5 = \qty{14.6}{g}$，约 \qty{15}{g}；
$m(\ce{CO2}) = 2.0 \times 44.0 = \qty{88}{g}$.
\end{solution}

\begin{solution}{exo:g10:the-mole:4}
$N = 0.50 \times \num{6.02e23} = \num{3.0e23}$ 个氧气分子。
$n = \num{3.0e24} / \qty{6.02e23}{mol^{-1}} = \qty{5.0}{mol}$，即铁的物质的量。
\end{solution}

\begin{solution}{exo:g10:the-mole:5}
$V = n \times V_m = 0.50 \times 24.1 = \qty{12}{L}$.
$n = V / V_m = 12 / 24.1 = \qty{0.50}{mol}$.
\end{solution}

\begin{solution}{exo:g10:the-mole:6}
液滴的质量为 \qty{0.050}{g}，故
$n = 0.050 / 18.0 = \qty{2.8e-3}{mol}$，且
$N = \num{2.8e-3} \times \num{6.02e23} = \num{1.7e21}$ 个分子。
\end{solution}

\begin{solution}{exo:g10:the-mole:7}
一个铝原子（\qty{27.0}{g/mol}）的质量约为
一个铁原子（\qty{55.8}{g/mol}）的一半，故等质量的铝所含
原子数约为铁的两倍。核对：铝 $10/27.0 = \qty{0.370}{mol}$，即
\num{2.23e23} 个原子；铁 $10/55.8 = \qty{0.179}{mol}$，即
\num{1.08e23} 个原子。铝的原子数约为铁的 2.1 倍。
\end{solution}

\begin{solution}{exo:g10:the-mole:8}
一升气体的物质的量为 $1/24.1 = \qty{0.0415}{mol}$。二氧化碳：
$0.0415 \times 44.0 = \qty{1.83}{g}$；氮气：
$0.0415 \times 28.0 = \qty{1.16}{g}$。二氧化碳的密度约为空气的一倍半，
而空气的摩尔质量与氮气相近。
发酵的果汁会释放出二氧化碳，它会下沉
并在封闭地窖近地面处积聚，可能使
进入其中的人窒息。
\end{solution}

\begin{solution}{exo:g10:the-mole:9}
$M(\ce{C6H12O6}) = 6 \times 12.0 + 12 \times 1.0 + 6 \times 16.0 =
\qty{180.0}{g/mol}$。箭头 $\div N_A$：$n = \num{1.0e22} / \num{6.02e23} =
\qty{1.66e-2}{mol}$；箭头 $\times M$：$m = \num{1.66e-2} \times 180.0 =
\qty{3.0}{g}$.
\end{solution}

\begin{solution}{exo:g10:the-mole:10}
$M = \qty{342.0}{g/mol}$，故 $n = 5.0 / 342.0 = \qty{1.46e-2}{mol}$；
$N = \num{1.46e-2} \times \num{6.02e23} = \num{8.8e21}$ 个分子。每个
含 12 个碳原子：$12 \times \num{8.8e21} = \num{1.1e23}$ 个碳原子。
\end{solution}

\begin{solution}{exo:g10:the-mole:11}
水：$1000 / 18.0 = \qty{55.6}{mol}$。乙醇：
$M = 2 \times 12.0 + 6 \times 1.0 + 16.0 = \qty{46.0}{g/mol}$，故
$1000 / 46.0 = \qty{21.7}{mol}$。水的物质的量更大，比值为
$55.6 / 21.7 \approx 2.6$，恰好等于 $46.0 / 18.0$。
\end{solution}

\begin{solution}{exo:g10:the-mole:12}
金：$\qty{3.75}{g}$，$n = 3.75 / 197.0 = \qty{1.90e-2}{mol}$，即
\num{1.15e22} 个原子。铜：\qty{1.25}{g}，
$n = 1.25 / 63.5 = \qty{1.97e-2}{mol}$，即 \num{1.19e22} 个原子。
出人意料的是，这枚戒指所含铜原子略多于金原子：
一个金原子的质量约为一个铜原子的三倍。
\end{solution}

\begin{solution}{exo:g10:the-mole:13}
\begin{enumerate}
  \item $V = 8.0 \times 6.0 \times 3.0 = \qty{144}{m^3} =
        \qty{1.44e5}{L}$; $n = \num{1.44e5} / 24.1 = \qty{5.98e3}{mol}$;
        $N = \num{5.98e3} \times \num{6.02e23} = \num{3.6e27}$ 个分子。
  \item $M = (78 \times 28.0 + 21 \times 32.0 + 1 \times 40.0) / 100 =
        2896 / 100 = \qty{29.0}{g/mol}$.
  \item $m = \num{5.98e3} \times 29.0 = \qty{1.73e5}{g}$，约
        \qty{173}{kg}：一间教室的空气重量与
        两三个成年人相当。
\end{enumerate}
\end{solution}

\begin{solution}{exo:g10:the-mole:14}
\begin{enumerate}
  \item $n = 1000 / 28.1 = \qty{35.6}{mol}$，故
        $N = 35.6 \times \num{6.02e23} = \num{2.14e25}$ 个原子。
  \item $m = \qty{1.000}{kg} / \num{2.14e25} = \qty{4.67e-26}{kg}$.
  \item 球的质量与摩尔质量可给出其物质的量，
        $n = m / M$。若原子数 $N$ 另行独立求出（由
        球的体积及晶体中原子的间距推算），
        则比值 $N / n$ 即为阿伏伽德罗常数。
\end{enumerate}
\end{solution}

\begin{solution}{exo:g10:the-mole:15}
$M = 0.758 \times 35.0 + 0.242 \times 37.0 = 26.53 + 8.95 =
\qty{35.5}{g/mol}$：即表中的数值。
\end{solution}

\begin{solution}{pb:g10:the-mole:1}
\textbf{1.} $M(\ce{H2O}) = 2 \times 1.0 + 16.0 = \qty{18.0}{g/mol}$.

\textbf{2.} $m = 250 \times \qty{1.0}{g} = \qty{250}{g}$.

\textbf{3.} $n = 250 / 18.0 = \qty{13.9}{mol}$.

\textbf{4.} 每个分子含两个氢原子和一个氧原子：
\qty{27.8}{mol} 氢原子，\qty{13.9}{mol} 氧原子。

\textbf{5.} $N = 13.9 \times \num{6.02e23} = \num{8.36e24}$ 个分子。

\textbf{6.} $m = \qty{18.0}{g/mol} / \qty{6.02e23}{mol^{-1}} =
\qty{2.99e-23}{g}$.

\textbf{7.} $\num{8.36e24} \times \qty{2.99e-23}{g} = \qty{250}{g}$，这正是
玻璃杯中水的质量，与预期相符。

\textbf{8.} $\num{8.36e24} / \num{3.16e7} = \num{2.6e17}$ 年：一个
含十七个零的年数，远超任何可能的计数。

\textbf{9.} $\qty{1.335e9}{km^3} \times \num{e9} = \qty{1.335e18}{m^3}$,
已知 $\qty{1}{m^3} = \qty{1000}{L}$，故 \qty{1.335e21}{L}。

\textbf{10.} 每升标记分子数为 $\num{8.36e24} / \num{1.335e21} = \num{6.26e3}$
个。

\textbf{11.} 第二杯中标记分子数为 $\num{6.26e3} \times 0.250 = \num{1.57e3}$
个。

\textbf{12.} 质量 $\num{1.335e21} \times \qty{1.0}{kg} =
\qty{1.335e21}{kg} = \qty{1.335e24}{g}$; $n = \num{1.335e24} / 18.0 =
\qty{7.42e22}{mol}$.

\textbf{13.} $13.9 / \num{7.42e22} = \num{1.87e-22}$.

\textbf{14.} $\num{1.87e-22} \times \num{8.36e24} = \num{1.57e3}$：
结果与第 11 题相同，这是必然的。

\textbf{15.} $\num{1.335e21} / 0.250 = \num{5.34e21}$ 杯。

\textbf{16.} $\num{8.36e24} / \num{5.34e21} \approx 1570$。一杯水
所含分子数约为海洋所含杯数的 1600 倍；
若分子均匀散布，则这一杯水的分子
在每满杯海水中约留有 1600 个。

\textbf{17.} $1 / \num{6.26e3} = \qty{1.6e-4}{L} = \qty{0.16}{mL}$，约
三滴。

\textbf{18.} 海洋体积只是估算值，至多只有三四位有效数字，且经过了取整；
海水并非纯水（其中溶有盐类）；
玻璃杯的容积并不恰好为 \qty{250}{mL}；
且混合在整个海洋中也不可能完全均匀。只有
数量级是可信的。

\textbf{19.} 第一杯中约 1600 个分子
会重新回到第二杯中。
\end{solution}
